%% ── ХАВСРАЛТ ─────────────────────────────────────────────────
\appendix
\chapter{Хавсралт}
\label{ch:appendix}

\section{Монгол хэл дэх хураангуй}

\noindent
Энэхүү судалгааны ажил нь хиймэл оюун ухааны технологийг ашиглан хэрэглэгчийн
тэмдэглэлийн агуулгад тохирсон автомат стикер үүсгэж, тэдгээр стикерт тулгуурласан
сануулга болон мэдэгдлийн системийг бүрдүүлсэн дижитал ухаалаг тэмдэглэлийн дэвтрийн
аппликейшнийг боловсруулахад чиглэсэн.

Системийн backend хэсгийг \textbf{Node.js + Express.js} ашиглан хөгжүүлж,
\textbf{MongoDB Atlas} үүлэн өгөгдлийн санд холбосон. Хэрэглэгчийн баталгаажуулалтыг
\textbf{JWT} болон \textbf{bcrypt} технологиор хэрэгжүүлсэн. Frontend хэсгийг
\textbf{Flutter (Dart)} хэлээр бичиж, Android болон iOS платформд ажиллах мобайл
аппликейшн бүтээв.

Судалгааны ажлын гол оруулт нь:
\begin{enumerate}
  \item Хэрэглэгчийн mood (сэтгэлийн байдал)-д тохирсон тэмдэглэлийн систем
  \item JWT-д суурилсан аюулгүй баталгаажуулалт
  \item REST API архитектур дээр суурилсан client-server систем
  \item Flutter-ийн widget tree ашиглан хялбар UI/UX
\end{enumerate}

Туршилтын үр дүнгээр API хариу хугацаа дундажаар 91 мс, хэрэглэгчийн сэтгэл ханамж
4.5/5.0 (90\%) байсан.

\vspace{1cm}
\section{Abstract in English}

\noindent
This research focuses on developing a digital smart note-taking mobile application that
automatically generates AI-powered stickers based on note content, and delivers reminder
notifications utilizing those stickers.

The backend was built with \textbf{Node.js + Express.js} connected to \textbf{MongoDB Atlas}
cloud database. User authentication was implemented using \textbf{JWT} and \textbf{bcrypt}.
The frontend was developed in \textbf{Flutter (Dart)}, supporting both Android and iOS
platforms.

Key contributions of this research:
\begin{enumerate}
  \item Mood-aware note-taking system that links notes to emotional states
  \item Secure JWT-based authentication with bcrypt password hashing
  \item Client-server architecture based on REST API principles
  \item Intuitive UI/UX using Flutter's widget tree with Material Design
\end{enumerate}

Testing results showed an average API response time of 91 ms and a user satisfaction
score of 4.5/5.0 (90\%).

\vspace{1cm}
\section{Бүтцийн диаграм -- Фолдерийн зохион байгуулалт}

\textbf{ai-notebook-backend фолдер:}
\begin{lstlisting}
ai-notebook-backend/
  server.js          <- Express сервер, MongoDB холболт
  package.json
  config/
  controllers/
    authController.js
    noteController.js  <- CRUD функцүүд
  middleware/
    authMiddleware.js  <- JWT шалгах
  models/
    Note.js           <- Mongoose schema
    User.js           <- Mongoose schema
  routes/
    auth.js           <- /api/auth endpoints
    notes.js          <- /api/notes endpoints
\end{lstlisting}

\textbf{ai\_notebook\_app фолдер (Flutter):}
\begin{lstlisting}
ai_notebook_app/
  lib/
    main.dart          <- App entry, SplashRouter
    login_screen.dart  <- Нэвтрэх UI
    Signup_screen.dart <- Бүртгэл UI
    Dashboard_screen.dart <- Тэмдэглэлийн жагсаалт
    main_screen.dart   <- Navigation
    services/
      api_service.dart <- HTTP requests, token storage
  pubspec.yaml
  android/
  ios/
\end{lstlisting}

\section{authMiddleware.js -- JWT шалгах middleware}

\begin{lstlisting}[caption={middleware/authMiddleware.js}]
const jwt = require("jsonwebtoken");
const JWT_SECRET = "YOUR_SECRET_KEY";

module.exports = (req, res, next) => {
  const authHeader = req.headers.authorization;
  if (!authHeader)
    return res.status(401).json({ message: "No token" });

  const token = authHeader.split(" ")[1];
  try {
    const decoded = jwt.verify(token, JWT_SECRET);
    req.userId = decoded.userId;
    next();
  } catch (err) {
    return res.status(401).json({ message: "Invalid token" });
  }
};
\end{lstlisting}
